\chapter{System Architecture and Problem Formulation}

% ==============================================================
\section{Notation}
% ==============================================================

\begin{table}[H]
\centering
\caption{Mathematical notation used throughout the report}
\begin{tabular}{p{2.8cm} p{10cm}}
\toprule
\textbf{Symbol} & \textbf{Definition} \\
\midrule
$\mathcal{G} = (V, E)$
  & Directed road graph; $V$ = intersections, $E$ = road segments \\
$e \in E$
  & A directed road edge (segment) \\
$\mathcal{A} = \{a_1, \dots, a_N\}$
  & Set of $N$ ingested articles / traffic incidents \\
$\mathcal{D} = \{d_1, \dots, d_M\}$
  & Set of $M$ extracted disruption events after LLM processing \\
$\sigma_i \in \{2, 5, 10\}$
  & Severity score of event $d_i$: low $= 2$, medium $= 5$, high $= 10$ \\
$\kappa_i \in [0, 1]$
  & Confidence weight of event $d_i$ \\
$\hat{\kappa}_i$
  & HGNN-adjusted confidence for event $d_i$ \\
$\kappa_i^{\text{final}}$
  & Blended confidence after HGNN correction \\
$Z_e(t) \in \{0, 1\}$
  & Binary disruption state of edge $e$ at time $t$ \\
$\pi_e^{\text{prior}}(t)$
  & Time-of-day prior disruption probability for edge $e$ \\
$\pi_e^{\text{post}}(t)$
  & Posterior disruption probability after Bayesian update \\
$\pi_e^{\text{final}}(t)$
  & Final per-edge probability after HGNN blend \\
$\mathcal{R} = \{r_1, \dots, r_K\}$
  & Set of $K$ candidate routes computed for a query \\
$\rho(r)$
  & Aggregate risk score of route $r$ \\
$t_r$
  & Expected travel time of route $r$ \\
$\mathcal{C}(r)$
  & Composite route cost used for ranking \\
$d_m, r_m$
  & Duration multiplier and recency multiplier \\
$\mathbf{h}_v^{(\phi)}$
  & Hidden embedding of node $v$ of type $\phi$ in the HGNN \\
$\alpha_{uv}$
  & Attention coefficient between nodes $u$ and $v$ \\
$\beta_\psi$
  & Semantic attention weight for meta-path $\psi$ \\
\bottomrule
\end{tabular}
\end{table}

% ==============================================================
\section{Formal Problem Definition}
% ==============================================================

Let $\mathcal{G} = (V, E)$ be the directed road graph of Kolkata.
Each edge $e \in E$ carries a free-flow travel time $t_e^0$ and a
disruption state $Z_e(t) \in \{0, 1\}$, which is latent and
must be inferred from external signals.

\medskip
\noindent
A user query $q = (\text{origin}, \text{destination}, \text{mode}, t)$
arrives at time $t$. The system must return a ranked list of routes
$r_1 \prec r_2 \prec \cdots \prec r_K$ from origin to destination
under mode constraints, where the ranking criterion is a composite
cost that accounts for both expected travel time and disruption risk.

\medskip
\noindent
The key challenge is that $Z_e(t)$ is not directly observable.
It must be estimated from unstructured signals — news articles,
RSS feeds, official advisories — that carry noisy, partial, and
sometimes conflicting information about road conditions.

\medskip
\noindent
\textbf{Objective.} Find the route $r^*$ that minimises the
risk-adjusted composite cost:
\begin{equation}
  r^* \;=\; \arg\min_{r \in \mathcal{R}}\; \mathcal{C}(r)
  \;=\; \arg\min_{r \in \mathcal{R}}\;
        t_r \cdot \Bigl(1 + \frac{\rho(r)}{10}\Bigr)
  \label{eq:objective}
\end{equation}

where $\rho(r)$ is the aggregate risk score of route $r$, derived from
the posterior edge disruption probabilities $\pi_e^{\text{final}}(t)$
computed through the three-layer pipeline described below.

% ==============================================================
\section{Layer Overview}
% ==============================================================

The system processes each query through three sequential layers.

\begin{table}[H]
\centering
\caption{System layer summary}
\begin{tabular}{clp{5cm}c}
\toprule
\textbf{Layer} & \textbf{Function} & \textbf{Key output} & \textbf{Status} \\
\midrule
1 & LLM extraction + HGNN refinement
  & $(\sigma_i, \kappa_i^{\text{final}})$ per event $d_i$
  & Complete \\
2 & Bayesian probabilistic fusion
  & $\pi_e^{\text{final}}(t)$ per edge $e$
  & Complete \\
3 & CVaR risk-aware routing
  & ranked routes $r_1 \prec \cdots \prec r_K$
  & In progress \\
\bottomrule
\end{tabular}
\end{table}

\noindent
The end-to-end pipeline for a single user query is:

\begin{equation}
  \mathcal{A}
  \;\xrightarrow{\text{pre-filter}}\;
  \tilde{\mathcal{A}}
  \;\xrightarrow{\text{LLM}}\;
  \mathcal{D}
  \;\xrightarrow{\text{HGNN}}\;
  \{\kappa_i^{\text{final}}\}
  \;\xrightarrow{\text{Bayes}}\;
  \{\pi_e^{\text{final}}\}
  \;\xrightarrow{\text{score}}\;
  \{\rho(r)\}
  \;\xrightarrow{\text{rank}}\;
  r^*
  \label{eq:pipeline}
\end{equation}

% ==============================================================
\section{Data Ingestion Sources}
% ==============================================================

\begin{longtable}{p{4.0cm} p{2.6cm} p{6.4cm}}
\toprule
\textbf{Source} & \textbf{Type} & \textbf{What it provides} \\
\midrule
\endhead
TomTom Incidents API v5~\cite{tomtom2024}
  & Structured, live
  & Accidents, closures, road-works with GPS coordinates;
    2{,}500 free req/day \\[4pt]
Google News RSS \texttt{when:2d}
  & Unstructured
  & City-wide Kolkata disruption news scoped to last 2 days \\[4pt]
NewsAPI
  & Unstructured
  & Per-road keyword search, last 24 h \\[4pt]
OpenWeatherMap~\cite{openweathermap2024}
  & Structured, live
  & Precipitation intensity, wind speed, visibility \\[4pt]
Kolkata Police (via RSS)
  & Unstructured
  & Official traffic advisories \\[4pt]
KMRC~\cite{kmrc2024}
  & Unstructured
  & Metro service disruption notices \\[4pt]
WB Disaster Management
  & Unstructured
  & Flood and cyclone alerts \\[4pt]
Indian Railways (erail.in)
  & Structured
  & Delayed trains at Howrah / Sealdah \\[4pt]
KMC Waterlogging (RSS)
  & Unstructured
  & Waterlogging hotspot reports \\
\bottomrule
\end{longtable}

% ==============================================================
\section{Two-Step API Design}
% ==============================================================

Equation~\eqref{eq:pipeline} has two distinct latency regimes.
Route geometry computation on the cached OSM graph takes
$\sim$2~s; the LLM pipeline over all articles takes 30--60~s.
Blocking the user interface on the slower step is unacceptable,
so the pipeline is split at the HGNN boundary:

\begin{table}[H]
\centering
\caption{Two-step decoupled request flow}
\begin{tabular}{p{0.46\textwidth} p{0.46\textwidth}}
\toprule
\textbf{Step 1 --- \texttt{POST /api/route} ($\sim$2 s)} &
\textbf{Step 2 --- \texttt{POST /api/disruptions} ($\sim$30--60 s)} \\
\midrule
Computes $\mathcal{R}$ using OSMnx Yen's $k$-shortest paths.
Returns GeoJSON immediately.
All routes assigned placeholder $\rho(r) = 0$, shown green. &
Executes steps $\tilde{\mathcal{A}} \to \mathcal{D} \to \{\kappa_i^{\text{final}}\}
\to \{\pi_e^{\text{final}}\} \to \{\rho(r)\}$.
Map updated with risk colours and disruption markers. \\
\bottomrule
\end{tabular}
\end{table}

% ==============================================================
\section{Transport Modes and Graph Construction}
% ==============================================================

For each transport mode $m$, a separate graph $\mathcal{G}_m = (V_m, E_m)$
is maintained:

\begin{table}[H]
\centering
\caption{Graph structure per transport mode}
\begin{tabular}{l l l}
\toprule
\textbf{Mode} & \textbf{Graph} & \textbf{Edge weight $w_e$} \\
\midrule
\texttt{drive}        & OSMnx drive graph        & Free-flow travel time \\
\texttt{walk}         & OSMnx pedestrian graph   & Walking time at 5 km/h \\
\texttt{bike}         & OSMnx cycling graph      & Cycling time at 15 km/h \\
\texttt{metro}        & 5-line station graph     & Scheduled travel time + headway \\
\texttt{metro+walk}   & Drive + metro + walk     & Composite (access + ride + egress) \\
\texttt{metro+bike}   & Bike + metro + bike      & Composite \\
\texttt{metro+drive}  & Drive + metro + drive    & Composite \\
\texttt{bus}          & Bus stop graph           & Estimated inter-stop time \\
\bottomrule
\end{tabular}
\end{table}

For multimodal journeys the total composite cost is:
\begin{equation}
  t_r^{\text{multimodal}}
  \;=\; t_r^{\text{access}} + t_r^{\text{transit}} + t_r^{\text{egress}}
  \label{eq:multimodal}
\end{equation}

where each leg uses the edge weights appropriate to its mode.

% ==============================================================
\section{Complete API Surface}
% ==============================================================

\begin{table}[H]
\centering
\caption{REST API endpoints --- built with FastAPI~\cite{fastapi2024}}
\begin{tabular}{l l p{7cm}}
\toprule
\textbf{Method} & \textbf{Endpoint} & \textbf{Description} \\
\midrule
GET  & \texttt{/}                            & Leaflet.js~\cite{leaflet2024} map frontend \\
GET  & \texttt{/api/locations}               & Localities, metro stations, bus stops \\
POST & \texttt{/api/route}                   & Step 1 --- compute $\mathcal{R}$ only \\
POST & \texttt{/api/disruptions}             & Step 2 --- full pipeline, returns $\rho(r)$ \\
POST & \texttt{/api/explain-route}           & HGNN attention weights per event \\
GET  & \texttt{/api/metro-overlay}           & GeoJSON of all 5 metro lines \\
GET  & \texttt{/api/metro-lines}             & Timetable summary per line \\
GET  & \texttt{/api/next-metro/\{station\}}  & Next $n$ trains (live IST time) \\
GET  & \texttt{/api/bus-overlay}             & GeoJSON of all bus stops \\
GET  & \texttt{/api/bus-network}             & Full bus route and stop data \\
GET  & \texttt{/api/hgnn-status}             & Model weights loaded / ready flag \\
\bottomrule
\end{tabular}
\end{table}
